\documentclass{article}
\usepackage[utf8]{inputenc}
\usepackage[T1]{fontenc}
\usepackage{amsfonts,amsmath,amssymb,amsthm}
\usepackage{hyperref}
\usepackage{booktabs}
\usepackage{graphicx}

\newcommand{\Var}{\mathrm{Var}}
\newcommand{\SHAP}{\mathrm{SHAP}}
\newcommand{\DASH}{\textsc{Dash}}

\newtheorem{theorem}{Theorem}
\newtheorem{lemma}[theorem]{Lemma}
\newtheorem{proposition}[theorem]{Proposition}
\newtheorem{axiom}{Axiom}

\title{Supplementary Material:\\The Attribution Impossibility}
\author{}
\date{}

\begin{document}
\maketitle

\section{Complete Axiom System}

Table~\ref{tab:axioms} lists the 6 property axioms in the formalization
(the core impossibility theorem requires none of them---only the Rashomon
property as hypothesis).
The formalization uses 12 axioms total: 5 type/constant declarations
(\texttt{Model}, \texttt{numTrees}, \texttt{numTrees\_pos}, \texttt{attribution},
\texttt{splitCount}, \texttt{firstMover}), the 5 property axioms below,
the DASH symmetry axiom, and one quantitative Spearman bound.

\begin{table}[h]
\centering
\caption{Complete axiom inventory.}
\label{tab:axioms}
\begin{tabular}{@{}llp{6cm}@{}}
\toprule
\textbf{Axiom} & \textbf{Lean~4 name} & \textbf{Justification} \\
\midrule
First-mover surjectivity & \texttt{firstMover\_surjective} & DGP symmetry: each feature can be first-mover \\
Split count (first-mover) & \texttt{splitCount\_firstMover} & Gaussian conditioning (Lemma~1 in main text) \\
Split count (non-first-mover) & \texttt{splitCount\_nonFirstMover} & Residual signal after first-mover absorbs $\rho$-component \\
Proportionality & \texttt{attribution\_proportional} & Uniform contribution model (Assumption~7) \\
Spearman bound & \texttt{spearman\_classical\_bound} & Combinatorial rank argument: tied non-first-movers \\
Attribution sum symmetry & \texttt{attribution\_sum\_symmetric} & DGP symmetry + balanced ensemble hypothesis \\
\bottomrule
\end{tabular}
\end{table}

\section{Proof Architecture}

The formalization comprises 13 Lean~4 files (12 axioms, 28 substantive theorems):

\begin{verbatim}
Defs.lean (11 axioms, type definitions)
  |
  +-- Trilemma.lean (RashimonProperty, attribution_impossibility)
  |     |
  |     +-- Iterative.lean (IterativeOptimizer, iterative_impossibility)
  |     |     +-- Lasso.lean (lasso_impossibility)
  |     |     +-- NeuralNet.lean (nn_impossibility)
  |     |
  |     +-- General.lean (GBDT instance, gbdt_impossibility)
  |
  +-- SplitGap.lean (split_gap_exact, split_gap_ge_half)
  +-- Ratio.lean (attribution_ratio, ratio_tendsto_atTop)
  +-- SpearmanDef.lean (spearmanCorr definition, derived bound)
  |
  +-- Impossibility.lean (not_equitable, not_stable, impossibility)
  +-- Corollary.lean (consensus_equity, consensus_difference_zero)
  +-- RandomForest.lean (contrast case, documentation only)
\end{verbatim}

\section{Inconsistencies Found During Formalization}

Lean's type checker caught three issues in the initial axiom system:

\begin{enumerate}
\item \textbf{First-mover balance (Axiom~6, original):}
  Originally stated universally over all model arrays---a constant
  model function trivially derived $\bot$ (the first-mover count for
  two different features cannot both equal the total). \emph{Fix:}
  Replaced with an explicit \texttt{IsBalanced} predicate as a
  hypothesis.

\item \textbf{Attribution sum symmetry (Axiom~7, original):}
  Combined with Axioms~2--4, the original unconditional version
  derived $\bot$ for unbalanced ensembles: the summed attributions
  cannot be equal when the first-mover counts differ. \emph{Fix:}
  Conditioned on \texttt{IsBalanced}.

\item \textbf{Split count type:}
  Originally \texttt{splitCount} returned $\mathbb{N}$, but
  $T/(2-\rho^2)$ is generally irrational, creating a type mismatch.
  \emph{Fix:} Changed to $\mathbb{R}$ (idealized leading-order values).
\end{enumerate}

The first two are genuine logical inconsistencies that would have been
difficult to detect by inspection of an informal proof. The third is a
modeling assumption made explicit by the type system.

\section{Gaussian Conditioning Argument}

The split count formulas (Axioms~2--3) are justified by the following
argument. Consider two features $X_j, X_k$ with correlation $\rho$.
At tree $t$, the available signal for feature $k$ after feature $j$
has been selected as root split is:
\[
  X_k \mid X_j = X_k - \rho X_j,
\]
which has variance $1 - \rho^2$. The expected number of splits for the
first-mover sums a geometric series over $T$ trees, yielding
$T/(2 - \rho^2)$. The non-first-mover receives the residual signal,
yielding $(1-\rho^2)T/(2-\rho^2)$. Full derivation with explicit
series computation is in the companion paper.

\section{Proof Sketch: $\alpha = 2/\pi$ from Gaussian Binary Quantization}

The main text reports $\alpha \approx 2/\pi$ for the signal capture fraction
of gradient-boosted stumps. This value has a clean theoretical derivation
from quantization theory.

\begin{proposition}
For $X \sim \mathcal{N}(0, \sigma^2)$, the optimal binary quantizer
(split at $x = 0$) captures fraction $2/\pi$ of the variance:
$\Var(\hat{X}) = (2/\pi)\sigma^2$.
\end{proposition}

\begin{proof}
The optimal two-level quantizer of a symmetric distribution splits at the
median ($x = 0$ for Gaussians). The quantized signal is:
$\hat{X} = \mathbb{E}[X \mid X > 0] \cdot \mathbf{1}_{X > 0} +
\mathbb{E}[X \mid X \leq 0] \cdot \mathbf{1}_{X \leq 0}$.
By symmetry, $\mathbb{E}[X \mid X > 0] = \sigma\sqrt{2/\pi}$
(the mean of a half-normal distribution). Thus
$\hat{X}$ takes values $\pm\sigma\sqrt{2/\pi}$ each with probability $1/2$,
giving $\Var(\hat{X}) = (2/\pi)\sigma^2$.
The fraction of variance captured is $\Var(\hat{X})/\Var(X) = 2/\pi \approx 0.637$.
\end{proof}

\paragraph{Connection to the corrected ratio.}
In each boosting round, a stump makes one binary split on the selected
feature, capturing $2/\pi$ of that feature's remaining signal variance.
This gives the effective signal capture fraction $\alpha = 2/\pi$ in
the corrected ratio $1/(1-\alpha\rho^2)$.
The fitted value ($\alpha \approx 0.60$) is slightly below $2/\pi = 0.637$
because (i) stumps split on the empirical median rather than the
population median, and (ii) after the first boosting round, residuals
are no longer exactly Gaussian.

\section{Fisher Information and the Rashomon Property}

The Rashomon property has a classical statistics interpretation through
the Fisher information matrix.

\paragraph{Setup.}
Consider the linear model $Y = \beta_j X_j + \beta_k X_k + \varepsilon$
where $X_j, X_k$ are jointly Gaussian with $\Var(X_j) = \Var(X_k) = 1$
and $\text{Cov}(X_j, X_k) = \rho$. The Fisher information matrix for
$\boldsymbol{\beta} = (\beta_j, \beta_k)$ is:
\[
    I(\boldsymbol{\beta}) = \frac{1}{\sigma^2}
    \begin{pmatrix} 1 & \rho \\ \rho & 1 \end{pmatrix}
\]
with eigenvalues $\lambda_\pm = (1 \pm \rho)/\sigma^2$.

\paragraph{Near-singularity under collinearity.}
As $\rho \to 1$, the smaller eigenvalue $\lambda_- = (1-\rho)/\sigma^2 \to 0$.
The FIM becomes near-singular along the direction $\boldsymbol{v}_- = (1, -1)/\sqrt{2}$,
which corresponds to redistributing credit between $j$ and $k$.
The profile likelihood is flat along this direction: models assigning more
credit to $j$ vs.\ $k$ are statistically indistinguishable.

\paragraph{Connection to the Rashomon property.}
The flat likelihood ridge is precisely the Rashomon set---the set of
models achieving near-optimal loss that disagree on $\beta_j$ vs.\
$\beta_k$.
Any model class containing an $\varepsilon$-neighborhood of the MLE
necessarily contains models on both sides of the ridge: some with
$\hat\beta_j > \hat\beta_k$ and others with $\hat\beta_k > \hat\beta_j$.
This is the Rashomon property, derived from the classical
non-identifiability of regression coefficients under collinearity.

The Cram\'er--Rao bound
$\Var(\hat\beta_j - \hat\beta_k) \geq 1/\lambda_- = \sigma^2/(1-\rho)$
quantifies the attribution instability: as $\rho \to 1$, the variance of
the difference $\hat\beta_j - \hat\beta_k$ diverges, confirming that no
finite sample can reliably determine which feature is more important.

% proposition already defined in preamble
\newtheorem{definition}[theorem]{Definition}
\newtheorem{remark}[theorem]{Remark}

\section{Extended Results: Universality and Relaxation Characterization}

The results in this section extend the main paper's impossibility theorem
in two directions: (i) showing the Rashomon property is inevitable for
symmetric model classes, making the impossibility inescapable; and
(ii) partially characterizing the remaining two relaxation paths.

\subsection{Rashomon Inevitability Under Symmetry}

The Attribution Impossibility (Theorem~1 in the main text) holds for
any model class satisfying the Rashomon property. We now show this
property is inevitable for symmetric model classes.

\begin{definition}[Permutation closure]
\label{def:perm-closure}
A model class $\mathcal{F}$ is \emph{permutation-closed} within
group $\ell$ if for any $f \in \mathcal{F}$ and any permutation
$\pi$ of features within group $\ell$, the composed model
$f \circ \pi \in \mathcal{F}$.
\end{definition}

This condition is satisfied by any model class without built-in
feature ordering: neural networks (permute input neurons),
gradient-boosted trees (permute split features), Lasso (permute
covariates), and random forests. It fails only for model classes
with hard-coded feature preferences.

\begin{theorem}[Rashomon from symmetry]
\label{thm:rashomon-symmetry}
Let the DGP be symmetric within group $\ell$ (permuting features
$j \leftrightarrow k$ leaves the population loss invariant).
Let $\mathcal{F}$ be permutation-closed within group $\ell$.
Then for any model $f \in \mathcal{F}$ with
$\varphi_j(f) \neq \varphi_k(f)$, the permuted model
$f' = f \circ \pi_{jk}$ satisfies:
\begin{enumerate}
    \item $L(f') = L(f)$ (same population loss),
    \item $\varphi_k(f') = \varphi_j(f)$ and
          $\varphi_j(f') = \varphi_k(f)$ (attributions swap).
\end{enumerate}
In particular, if $\varphi_j(f) > \varphi_k(f)$, then
$\varphi_k(f') > \varphi_j(f')$, and the Rashomon property holds.
\end{theorem}

\begin{proof}
By DGP symmetry, the joint distribution of $(X, Y)$ is invariant
under permutation of features $j$ and $k$ within group $\ell$.
Therefore $L(f') = L(f \circ \pi_{jk}) = L(f)$, since the
population loss depends only on the joint distribution.
By construction, $f'$ applies $f$ to permuted inputs, so its
reliance on feature $k$ equals $f$'s reliance on feature $j$:
$\varphi_k(f') = \varphi_j(f)$ and vice versa.
\end{proof}

\begin{theorem}[Attribution non-degeneracy]
\label{thm:non-degeneracy}
Let $A$ be a stochastic training algorithm with continuous dependence
on its random seed (e.g., SGD with random initialization, bootstrap
sampling, random feature selection). For any DGP with $\rho > 0$ and
$n < \infty$ training samples, $\Pr[\varphi_j(f) \neq \varphi_k(f)] = 1$
for features $j, k$ in the same collinear group.
\end{theorem}

\begin{proof}
With finite samples, the empirical correlation between $X_j$ and $X_k$
is $\hat\rho = \rho + O(1/\sqrt{n})$, which is almost surely irrational.
Under continuous dependence on the data, the training algorithm's feature
utilization varies continuously with $\hat\rho$, so the event
$\{\varphi_j(f) = \varphi_k(f)\}$ has measure zero in the joint
randomness of the data and the algorithm.
\end{proof}

\begin{theorem}[Rashomon inevitability]
\label{thm:rashomon-inevitability}
Let $A$ be a stochastic, symmetric training algorithm
($A$ applied to a within-group-permuted dataset produces the
permuted model in distribution) for a permutation-closed model class
$\mathcal{F}$. For any DGP with $\rho > 0$, the Rashomon property
holds: for any $j, k$ in the same collinear group, there exist
models $f, f'$ in the support of $A$ with $\varphi_j(f) > \varphi_k(f)$
and $\varphi_k(f') > \varphi_j(f')$.
\end{theorem}

\begin{proof}
By Theorem~\ref{thm:non-degeneracy},
$\Pr[\varphi_j(f) \neq \varphi_k(f)] = 1$, so either
$\varphi_j > \varphi_k$ or $\varphi_k > \varphi_j$ almost surely.
By algorithmic symmetry,
$\Pr[\varphi_j > \varphi_k] = \Pr[\varphi_k > \varphi_j]$.
Since these probabilities sum to~1, each equals~$1/2$.
Both events have positive probability, so both are realized.
\end{proof}

\begin{remark}
Theorem~\ref{thm:rashomon-inevitability} makes the Attribution
Impossibility inescapable for any standard ML pipeline under
collinearity. The only escapes are: (i)~$\rho = 0$ (no collinearity),
(ii)~a deterministic model class producing a single model, or
(iii)~an asymmetric algorithm with built-in feature preferences.
\end{remark}

\subsection{The Price of Dropping Faithfulness}

The main text resolves the impossibility by dropping completeness
(DASH: ties for symmetric features). What if we instead drop
faithfulness---allowing a stable, complete ranking that does not
reflect any specific model?

\begin{theorem}[Unfaithfulness lower bound]
\label{thm:unfaithfulness}
Let $j, k$ be symmetric features (same group, same population-level
relationship with $Y$). Any stable, complete ranking $\succ$ has
expected unfaithfulness at least $1/2$ for the pair $(j,k)$:
\[
    \Pr_f\bigl[\, j \succ k \;\text{but}\;
    \varphi_k(f) > \varphi_j(f) \,\bigr] \;=\; \frac{1}{2},
\]
where the probability is over models $f$ drawn from the symmetric
training distribution.
\end{theorem}

\begin{proof}
By DGP symmetry, the distribution of $\varphi_j(f) - \varphi_k(f)$
is symmetric about zero:
$\Pr[\varphi_j(f) > \varphi_k(f)] = \Pr[\varphi_k(f) > \varphi_j(f)] = 1/2$.
Any stable ranking must fix $j \succ k$ or $k \succ j$ independently
of $f$. Whichever it chooses disagrees with exactly half the models.
\end{proof}

\begin{proposition}[Optimal unfaithful ranking]
\label{prop:optimal-unfaithful}
The ranking $j \succ^* k \Leftrightarrow \mathbb{E}[\varphi_j] > \mathbb{E}[\varphi_k]$
minimizes expected unfaithfulness over the symmetric model distribution.
For symmetric features, $\mathbb{E}[\varphi_j] = \mathbb{E}[\varphi_k]$,
so $\succ^*$ assigns a tie---recovering the ``drop completeness'' solution.
\end{proposition}

\begin{proof}
The ranking $\succ^*$ is the Bayes-optimal binary decision for
``is $\varphi_j(f) > \varphi_k(f)$?'' under the model distribution,
minimizing $0$-$1$ loss. For symmetric features the posterior is exactly
$1/2$, so the optimal decision is indeterminate (a tie).
\end{proof}

\begin{theorem}[Relaxation path convergence]
\label{thm:path-convergence}
The ``drop faithfulness'' and ``drop completeness'' relaxation paths
converge to the same solution: population-level attributions with ties
for symmetric features. Consequently, the impossibility trilemma
collapses to a single tradeoff axis parameterized by ensemble size~$M$:
\begin{itemize}
    \item At $M = 1$: faithfulness and completeness, but no stability;
    \item As $M \to \infty$: stability and between-group faithfulness,
          but within-group ties (completeness relaxed).
\end{itemize}
Between-group faithfulness is preserved at all~$M$.
\end{theorem}

\begin{proof}
By Proposition~\ref{prop:optimal-unfaithful}, the optimal stable
complete ranking assigns ties to symmetric features---identical to the
``drop completeness'' solution (DASH). Both paths yield population-level
attributions. The $M$-parameterization follows from DASH consensus:
$\bar\varphi_j = \frac{1}{M}\sum_{i=1}^M \varphi_j(f_i)$, which
interpolates between a single model ($M=1$, faithful but unstable)
and the population mean ($M \to \infty$, stable but tied within groups).
\end{proof}

\section{Depth Dependence: Full Experimental Results}

Table~\ref{tab:depth-rho} reports the within-group split count ratio
across tree depths and correlation levels (XGBoost, $T{=}100$, $\eta{=}1.0$,
$P{=}10$ in 2 groups of 5, $N{=}2000$, 30 seeds).

\begin{table}[h]
\centering
\caption{Within-group split count ratio by depth and $\rho$ ($\eta{=}1.0$).}
\label{tab:depth-rho}
\small
\begin{tabular}{@{}lcccccc@{}}
\toprule
 & $\rho{=}0.3$ & $\rho{=}0.5$ & $\rho{=}0.7$ & $\rho{=}0.9$ & $\rho{=}0.95$ & Fitted $\alpha$ \\
\midrule
Depth 1 (stumps) & 1.28 & 1.32 & 1.42 & 1.99 & 2.16 & $0.60\ (\approx 2/\pi)$ \\
Depth 3          & 1.19 & 1.26 & 1.24 & 1.30 & 1.32 & $0.30$ \\
Depth 6          & 1.65 & 1.67 & 1.74 & 1.90 & 2.03 & $0.60$ \\
Depth 10         & 2.23 & 2.25 & 2.28 & 2.49 & 2.62 & --- \\
\midrule
Theory ($\alpha{=}1$) & 1.10 & 1.33 & 1.96 & 5.26 & 10.26 & $1.00$ \\
\bottomrule
\end{tabular}
\end{table}

\paragraph{Depth 3 anomaly.}
At depth~3, each tree has up to 7 leaves and uses multiple features per
tree. Internal splits distribute split counts across group members,
counteracting the root split's first-mover advantage.
The ratio is nearly $\rho$-independent (${\approx}1.3$), explaining the
low fitted $\alpha = 0.30$.

\paragraph{Depth 10.}
At depth~10, the root split cascades: once the first-mover is selected
at the root, subsequent splits on the same feature are more likely at
deeper levels (residuals retain first-mover signal). The baseline ratio
(${\approx}2.2$ even at $\rho{=}0.3$) is high because deep trees
overfit to the first-mover's signal.
The $1/(1{-}\alpha\rho^2)$ model does not fit well for depth~10 (the
ratio has a large $\rho$-independent component); we omit the fitted $\alpha$.

\paragraph{Practitioner guidance.}
Depth~3 gives the \emph{lowest} attribution instability among standard
configurations. For users prioritizing explanation stability, shallow
ensembles (depth~3) with \DASH{} consensus ($M \geq 5$) provide the
best stability--accuracy tradeoff.

\section{Real-World Validation: Breast Cancer (Wisconsin)}

\begin{figure}[h]
\centering
\includegraphics[width=0.55\textwidth]{figures/real_world_instability.pdf}
\caption{SHAP ranking instability on the Breast Cancer (Wisconsin) dataset.
Each dot is one of the 435 feature pairs; the $x$-axis is the absolute
pairwise correlation. Highly correlated feature pairs with similar
true importance (e.g., worst perimeter $\leftrightarrow$ worst area,
$|\rho|{=}0.98$, flip rate 48\%) exhibit near-maximal instability,
while correlated pairs with distinct importance remain stable.
50 XGBoost models with independent random seeds.}
\label{fig:real-world}
\end{figure}

We validate the impossibility on the Breast Cancer (Wisconsin) dataset
(30 features, 569 samples), which exhibits natural collinearity among
measurements of cell nuclei (radius, perimeter, area at mean/SE/worst levels).
Figure~\ref{fig:real-world} plots the SHAP ranking flip rate for each
feature pair against their absolute correlation.
Pairs measuring the same underlying quantity---such as worst perimeter
and worst area ($|\rho| = 0.98$, flip rate 48\%)---approach the
theoretical maximum of 50\% predicted by our DGP symmetry argument.
Pairs with high correlation but distinct importance (e.g., mean area and
mean concave points, $|\rho| = 0.82$, flip rate 2\%) remain stable, confirming
that the impossibility requires both collinearity \emph{and} similar
true importance.
Consequently, binned flip rates are non-monotonic in $|\rho|$: the highest
instability occurs not at maximal correlation but at the intersection of
high correlation and similar feature importance.

\section{\DASH{} Convergence on Breast Cancer}

We demonstrate that \DASH{} resolves instability on real data.
Using the 50~XGBoost models trained for the Breast Cancer validation
(\S\ref{fig:real-world}), we measure the flip rate for the most
unstable pair (worst perimeter $\leftrightarrow$ worst area,
$|\rho|{=}0.978$) as a function of ensemble size~$M$.

\begin{center}
\small
\begin{tabular}{@{}lcccccc@{}}
\toprule
$M$ & 1 & 2 & 5 & 10 & 25 & 50 \\
\midrule
Flip rate & 0.448 & 0.466 & 0.438 & 0.452 & 0.298 & 0.000 \\
\bottomrule
\end{tabular}
\end{center}

At $M{=}1$ the flip rate is 44.8\%, near the 50\% theoretical maximum.
\DASH{} consensus at $M{=}50$ resolves to the true ordering (worst
perimeter slightly more important: mean $|\SHAP| = 0.940$ vs.\ $0.901$)
with zero flips.
The convergence is slower than in the synthetic setting (where
features are exactly symmetric) because here the features have a slight
importance asymmetry---\DASH{} correctly detects and resolves it rather
than producing a tie.

\section{Extended Experimental Details}

\begin{table}[h]
\centering
\caption{Experimental configurations.}
\small
\begin{tabular}{@{}lp{10cm}@{}}
\toprule
Experiment & Configuration \\
\midrule
Synthetic (main text) & $P{=}20$, $L{=}4$ groups of $m{=}5$,
  $N{=}2{,}000$, $T{=}100$, $\eta{=}0.1$, 50 seeds,
  $\rho \in \{0.1, 0.3, 0.5, 0.7, 0.9, 0.95\}$ \\
Validation (depth$\times\rho$) & $P{=}10$, $L{=}2$ groups of $m{=}5$,
  $N{=}2{,}000$, $T{=}100$, $\eta{=}1.0$, 30 seeds,
  depths $\{1,3,6,10\}$, $\rho \in \{0.3, 0.5, 0.7, 0.9, 0.95\}$ \\
Breast Cancer & 30 features, 569 samples, $T{=}100$,
  depth${=}6$, $\eta{=}0.1$, 50 seeds, TreeSHAP on 200 test samples \\
\midrule
Hardware & Apple Silicon (M-series), single core \\
Software & Python~3.9.6, XGBoost~2.1.4, SHAP~0.49.1, NumPy~2.0.2,
  scikit-learn \\
Compute & Synthetic: ${<}5$ min; Validation: ${<}30$ min;
  Breast Cancer: ${<}10$ min \\
\bottomrule
\end{tabular}
\end{table}

All experiments use publicly available datasets (California Housing and
Breast Cancer from scikit-learn) and are fully reproducible from the
scripts in the anonymized code repository.
A minimal reference implementation of \DASH{} consensus attribution
(\texttt{dash\_example.py}) is included in the code supplement,
demonstrating the practical recommendation of \S5: train $M \geq 25$
models, average SHAP values, and diagnose instability via ranking flips.
The full \DASH{} pipeline (population generation, $\varepsilon$-filtering,
deduplication) is available in the companion repository.

\section{SymPy Verification}

All algebraic consequences have been independently verified:

\begin{verbatim}
# dash-shap/paper/proofs/verify_lemma6_algebra.py
#
# Verifies:
#   split_gap = rho^2 * T / (2 - rho^2)
#   attribution_ratio = 1 / (1 - rho^2)
#   split_gap >= rho^2 * T / 2  (for rho in (0,1))
#
# Result: ALL CHECKS PASS
\end{verbatim}

\end{document}
